\chapter{Methodology}
\label{chap:methods}

\section{Experimental Environments}

The thesis experiments are primarily classified into two main agricultural domain environments: \textbf{Fertilization} and \textbf{Crop Planning}. 

\subsection{Fertilization Domain}
The fertilization simulations required agents to control the daily or weekly application of Nitrogen, Phosphorus, and Potassium (NPK) to crops. The objective was to balance the maximization of crop yield (deterministic return) with the minimization of fertilizer costs and environmental runoff. Environments were configured with distinct horizons, often set to 3000 to 5000 years, to assess long-term cyclic policy stability and handle extreme weather events in deep recurrences.

\subsection{Crop Planning Domain}
The crop planning task represented a higher-level, lower-frequency decision scheme, deciding what to plant each season to keep soil health optimal without relying purely on chemical intervention. A major innovation was the introduction of \textit{Hierarchical Guarded Reruns}, wherein high-level actions were constrained by explicit structural guardrails (compliance shaping) to prevent the agent from taking physically impossible or highly penalized actions.

\section{Algorithm Suite}

The reinforcement learning pipeline systematically evaluated:
\begin{itemize}
    \item \textbf{Proximal Policy Optimization (PPO)}: Analyzed for its robust clipping mechanism, generally providing stable updates. Evaluated in both adaptive and non-adaptive variants based on weather.
    \item \textbf{Advantage Actor-Critic (A2C)}: Deployed as a baseline on-policy actor-critic algorithm. Historically computationally cheaper, but evaluated heavily to see if it could match PPO across long horizons.
    \item \textbf{Deep Q-Network (DQN)}: Used as descriptive baseline reference points. Due to its off-policy discrete nature, it served primarily to anchor the lower bound, revealing the performance floor compared to advanced policy gradient methods.
\end{itemize}

\section{Stochasticity: Fixed vs. Random Weather}

A crucial dimension of the matrix was testing the algorithms against distinct environmental stochasticity profiles:
\begin{itemize}
    \item \textbf{Fixed Weather}: A repeatable, deterministic sequence of climate conditions, reducing variance to strictly algorithmic exploration and policy behavior.
    \item \textbf{Random Weather}: A highly stochastic environment injecting unpredictable rain, temperature drops, and drought scenarios, dramatically complicating the state transition $\mathcal{P}(s'|s, a)$.
\end{itemize}

\section{Evaluation Metrics and Validation Procedure}

The canonical metric utilized to rank the runs was the \textbf{Deterministic Return} and \textbf{Evaluation Reward}. Specifically, performance across grouped categories was distilled via paired t-tests or non-parametric alternatives when applicable. 
During the final extraction step, a strict quality assurance process validated completeness:
\begin{itemize}
    \item All recovered exports were matched with their corresponding \texttt{model.zip}.
    \item Vector normalization states (\texttt{vec\_normalize\_*.pkl}) were carefully persisted for inference.
    \item Metrics were reconstructed strictly from authoritative evaluation arrays logged parallel to checkpoints.
\end{itemize}
